\documentclass[13pt, a4paper]{extarticle}

\usepackage[utf8]{inputenc}
\usepackage[T1]{fontenc}
\usepackage{mathptmx}
\usepackage[english]{babel}
\usepackage{geometry}
\geometry{a4paper, left=3cm, right=2cm, top=2.5cm, bottom=2.5cm}
\usepackage{fancyhdr}
\pagestyle{fancy}
\fancyhf{}
\fancyhead[C]{\thepage}
\renewcommand{\headrulewidth}{0pt}
\usepackage{amsmath, amssymb}
\usepackage{graphicx}
\usepackage{caption}
\usepackage{booktabs}
\usepackage{indentfirst}
\usepackage[authoryear, round]{natbib}

\begin{document}

\begin{titlepage}
    \thispagestyle{empty}
    \begin{center}
        TON DUC THANG UNIVERSITY \\
        FACULTY OF INFORMATION TECHNOLOGY \\
        \vspace{2.2cm}

        Student Scientific Research Project 2025--2026 \\
        \vspace{1.4cm}

        \textbf{\Large A Plateau-Aware Deep RL Controller for ALNS on VRPTW (DDQN-ALNS)} \\
        \vspace{2cm}
    \end{center}

    \begin{flushleft}
        \textbf{Advisor:} Dr. Ho Thi Linh \\
        \vspace{0.8cm}
        \textbf{Students:} \\
        1. Huynh Nhat Huy \\
        2. Nguyen Nhat Huy \\
        3. Tran (to be updated)
    \end{flushleft}

    \vfill
    \begin{center}
        Ho Chi Minh City, April 2026
    \end{center}
\end{titlepage}

\newpage
\setcounter{page}{1}

\section*{Introduction}
The Vehicle Routing Problem with Time Windows (VRPTW) is a core optimization problem in logistics. Adaptive Large Neighborhood Search (ALNS) is effective in practice, but it often reaches \textit{plateau} phases where many iterations produce no best-solution improvement. In these phases, computation budget is consumed with limited value.

This work proposes \textit{PlateauHybridSolver} (DDQN-ALNS), where a Double DQN acts as a \textbf{phase-level controller}. The policy is activated only when ALNS shows stagnation (\texttt{no\_imp} $\ge 60$), and selects one of five search modes to rebalance exploration and exploitation.

\section{Literature Background}
The classical VRP formulation started with \citep{Dantzig1959}. Solomon benchmarks remain a standard for VRPTW evaluation \citep{Solomon1987}. Recent research combines deep RL and metaheuristics to improve operator control quality. In contrast to per-iteration operator selection, this work uses a lighter mode-level policy triggered only during plateau phases, reducing control overhead and improving transferability.

\section{Objectives and Method}
\subsection{Objectives}
\begin{enumerate}
    \item Reduce average distance gap (\%) versus BKS on Solomon RC1 and RC2.
    \item Reduce runtime under the same iteration budget.
    \item Evaluate zero-shot transfer: policy trained on RC1, then applied to RC2 without retraining.
\end{enumerate}

\subsection{DDQN-ALNS Architecture}
ALNS is split into segments (50 iterations each). After each segment, DDQN observes a 10-dimensional normalized state:
\[
S_t = \big[\text{patience}, \text{cost\_gap}, \text{temp\_ratio}, \text{imp\_rate}, \text{nv\_ratio},
\text{route\_balance}, \text{load\_balance}, \text{tw\_tight}, \text{avg\_slack}, \text{progress}\big]
\]
and chooses one mode from:
\begin{itemize}
    \item \textbf{Default}
    \item \textbf{Intensify}
    \item \textbf{Diversify}
    \item \textbf{TW Rescue}
    \item \textbf{Route Reduce}
\end{itemize}

\subsection{Reward Design}
Segment reward combines best-solution improvement, current-solution improvement, acceptance quality, and fleet-fill dynamics:
\[
R_t = -0.75 - 0.75 \cdot \mathbb{I}[\text{accepted}=0] + R^{best} + R^{cur} + 3(\rho_{after}-\rho_{before})
\]
where $R^{best}$ heavily prioritizes best-solution progress, and $R^{cur}$ adds smaller local shaping while penalizing unstable vehicle growth.

\subsection{Training Setup}
Two hidden layers (96 units), Adam ($3\times10^{-4}$), $\gamma=0.95$, and epsilon decay ($0.35\rightarrow0.05$). Replay buffer size is 4096, target network update every 20 steps.

\section{Results and Discussion}
All results are averaged over 5 independent runs per (instance, algorithm). The summary below matches the benchmark CSV used in the repository.

\begin{table}[h]
\centering
\caption{Aggregated results on Solomon RC1 and RC2.}
\label{tab:summary}
\begin{tabular}{llrrrrr}
\toprule
\textbf{Set} & \textbf{Algorithm} & \textbf{NV} & \textbf{TD} & \textbf{Gap\%} & \textbf{OnTime} & \textbf{Time} \\
\midrule
RC1 & ALNS & 12.5 $\pm$ 0.3 & 1393.5 $\pm$ 13.1 & +1.0\% & 100.0\% & 120.7s \\
RC1 & DDQN-ALNS & 12.7 $\pm$ 0.3 & \textbf{1387.3} $\pm$ 18.4 & \textbf{+0.5\%} & 100.0\% & \textbf{79.8s} \\
\midrule
RC2 & ALNS & 3.5 $\pm$ 0.1 & 1132.2 $\pm$ 23.8 & +1.6\% & 100.0\% & 91.3s \\
RC2 & DDQN-ALNS & 3.6 $\pm$ 0.1 & 1121.0 $\pm$ 32.3 & \textbf{+0.7\%} & 100.0\% & 64.4s \\
RC2 & DDQN-ALNS$^\star$ & 3.6 $\pm$ 0.1 & 1121.4 $\pm$ 31.9 & +0.7\% & 100.0\% & 64.2s \\
\bottomrule
\end{tabular}
\vspace{0.2em}
\small $^\star$ Zero-shot transfer policy trained on RC1 and applied to RC2.
\end{table}

\subsection{Key Findings}
\begin{itemize}
    \item \textbf{RC1 gap improvement is statistically significant}: Wilcoxon signed-rank test gives $p=0.0078$.
    \item \textbf{Runtime consistently improves}: about 30--34\% faster than ALNS in RC1/RC2.
    \item \textbf{Zero-shot transfer is stable}: DDQN-ALNS$^\star$ on RC2 is nearly identical to directly trained DDQN-ALNS.
    \item \textbf{Vehicle count improvement is limited}: NV does not significantly improve over ALNS in this setup.
\end{itemize}

\subsection{Limitations}
\begin{enumerate}
    \item Evaluation currently focuses on RC1 and RC2 only.
    \item NV reduction remains weak and requires additional reward/repair redesign.
    \item Distribution shift from Solomon to real city delivery data can reduce policy robustness.
\end{enumerate}

\section{Conclusion}
This project demonstrates that a phase-level DDQN controller can improve ALNS behavior during plateau phases in VRPTW. The strongest gains are runtime reduction and statistically meaningful distance-gap improvement on RC1, with promising zero-shot transfer to RC2. Future work should target stronger NV optimization, larger benchmark sets, and domain adaptation for real-world operations.

\newpage
\begin{thebibliography}{}

\bibitem[Dantzig and Ramser(1959)]{Dantzig1959}
Dantzig, G. B., and Ramser, J. H. (1959).
\textit{The Truck Dispatching Problem}.
Management Science, 6(1), 80--91.

\bibitem[Solomon(1987)]{Solomon1987}
Solomon, M. M. (1987).
\textit{Algorithms for the Vehicle Routing and Scheduling Problems with Time Window Constraints}.
Operations Research, 35(2), 254--265.

\bibitem[Chen et al.(2020)]{Chen2020Hybrid}
Chen, X., Ulmer, M. W., and Thomas, B. W. (2020).
\textit{Deep Q-Learning for Same-Day Delivery with Vehicles and Drones}.
European Journal of Operational Research.

\end{thebibliography}

\end{document}
